\section{Appraisal Entry System}

\subsection{Design}
The Appraisal Entry System facilitates the documentation, categorisation, and retrieval of generated clinical summaries. Following a consultation, the doctor must be able to review the generated summary, append clinical notes, assign metadata (such as categorical tags), and save the record for future reference. 

To satisfy both external compliance and internal application requirements, a dual-storage paradigm was designed:
\begin{itemize}
    \item \textbf{External Export}: A human-readable text file is generated and saved to a user-defined directory via the Windows File Explorer, maintaining the NHS standard \texttt{YYYY-MM-DD\_Appraisal\_} naming convention.
    \item \textbf{Internal State}: A structured JSON file is maintained within the application's isolated local storage environment. This acts as a lightweight database, enabling the application to load, filter, and modify historical appraisals across sessions without requiring a heavy embedded SQL database.
\end{itemize}

Furthermore, the system requires a robust search and filtering mechanism to allow practitioners to retrieve historical records by name or specific clinical tags, alongside the ability to modify or delete existing entries.

\subsection{Implementation}
The appraisal system is integrated into the \texttt{MainWindow} class structure and handled within the \texttt{MainWindow.Appraisal.cpp} and \texttt{MainWindow.DisplayAppraisals.cpp} translation units. It uses the \texttt{winrt::Windows::Data::Json} and \texttt{winrt::Windows::Storage} namespaces for data persistence, and dynamically constructs WinUI3 elements at runtime due to the static nature of a WinUI3 XAML front-end

\subsubsection{Data Serialisation and Storage}
When a user saves an appraisal, \texttt{SaveAppraisalToJsonAsync} manages the serialisation. The application retrieves \texttt{appraisals.json} from the MSIX \texttt{LocalFolder}. If the file is corrupted or missing, a new \texttt{JsonObject} is instantiated. 

A notable implementation detail is the processing of user-defined tags. The input string is parsed via a \texttt{std::wstringstream}, delimited by commas, and subjected to whitespace trimming before being appended to a \texttt{JsonArray}. If no tags are provided, the system defaults to an \texttt{untagged} label.

\begin{verbatim}
// string trimming and tokenisation for tags
while (std::getline(ts, token, L',')) {
    size_t first = token.find_first_not_of(L" ");
    size_t last = token.find_last_not_of(L" ");
    if (first != std::string::npos && last != std::string::npos) {
        token = token.substr(first, (last - first + 1));
        tagsArray.Append(JsonValue::CreateStringValue(token));
    }
}
// default to untagged if empty
if (tagsArray.Size() == 0) tagsArray.Append(JsonValue::CreateStringValue(L"untagged"));

appraisalData.SetNamedValue(L"tags", tagsArray);
rootObject.SetNamedValue(keyName, appraisalData);
co_await FileIO::WriteTextAsync(file, rootObject.Stringify());
\end{verbatim}

\subsubsection{In Memory Caching and Filtering}
Appraisals are loaded into a \texttt{std::vector<AppraisalItem> m\_allAppraisals} vector upon navigating to the history view via \texttt{LoadAppraisalsAsync()}.  Filtering is applied in-memory. The \texttt{RenderAppraisalsList()} method iterates over the cached vector, applying a substring match for the appraisal name and checking against a \texttt{std::set<std::wstring> m\_activeFilters} for tag matches. The resulting subset is sorted chronologically using \texttt{std::sort} prior to rendering.

\subsubsection{Dynamic User Interface Generation}
For each filtered \texttt{AppraisalItem}, the application iteratively constructs a \texttt{Button} containing a \texttt{Border} and a \texttt{StackPanel}. This approach allows for control over the glass style theme. Additionally, embedding the appraisal's unique identifier into the \texttt{Tag} property of the button allows click events to be handled by the generated button.

\subsubsection{State Synchronisation during Modification}
When an appraisal is edited or deleted, the system must synchronise the in-memory cache, the JSON storage, and the UI. For instance, within \texttt{HistoryDialog\_DeleteClick()}, the target entry is first removed from the \texttt{m\_allAppraisals}. Next, the JSON file is read, the target key is removed and the file is overwritten. Finally, \texttt{RenderAppraisalsList()} is invoked to rebuild the user interface from the updated cache, ensuring the deleted item is immediately removed from the view.

\begin{verbatim}
// remove from internal list
auto it = std::remove_if(m_allAppraisals.begin(), m_allAppraisals.end(),
    [&](const AppraisalItem& app) { return app.Name == targetName; });

if (it != m_allAppraisals.end()) {
    m_allAppraisals.erase(it, m_allAppraisals.end());
}

// remove from file on disk
IStorageItem item = co_await localFolder.TryGetItemAsync(L"appraisals.json");
StorageFile file = item.as<StorageFile>();
winrt::hstring jsonString = co_await FileIO::ReadTextAsync(file);
JsonObject rootObject;
if (JsonObject::TryParse(jsonString, rootObject)) {
    if (rootObject.HasKey(targetName)) {
        rootObject.Remove(targetName);
        co_await FileIO::WriteTextAsync(file, rootObject.Stringify());
    }
}
\end{verbatim}